% ---- pg-356 （书页 344） ----

\[
\bar{\mathbf{x}}'=[4.47\qquad3.96\qquad71.42\qquad26.91\qquad1.64]
\]

\begin{center}
\begin{tabular}{ccccc}
总人口 & 职业等级 & 16 岁以上 & 政府雇佣 & 家庭\\
(千) & (百分比) & 人员的就业率 & (百分比) & 收入中位数(10 万元)\\
& & (百分比) & &
\end{tabular}
\end{center}

以及
\[
\mathbf{S}=\begin{bmatrix}
3.397&-1.102&4.306&-2.078&0.027\\
-1.102&9.673&-1.513&10.953&1.203\\
4.306&-1.513&55.626&-28.937&-0.044\\
-2.078&10.953&-28.937&89.067&0.957\\
0.027&1.203&-0.044&0.957&0.319
\end{bmatrix}
\]
样本变差能否用一或二个主成分来综合? 我们得到以下结果:

\begin{mvtable}
\begin{center}\mvbold{主成分系数(括号中是相关系数)}\end{center}
\begin{tabular}{lccccc}
\toprule
变量 & $\hat{\mathbf{e}}_1(r_{\hat{y}_1,x_k})$ & $\hat{\mathbf{e}}_2(r_{\hat{y}_2,x_k})$ & $\hat{\mathbf{e}}_3$ & $\hat{\mathbf{e}}_4$ & $\hat{\mathbf{e}}_5$\\
\midrule
总人口 & $-0.039(-0.22)$ & $0.071(0.24)$ & $0.188$ & $0.977$ & $-0.058$\\
职业 & $0.105(0.35)$ & $0.130(0.26)$ & $-0.961$ & $0.171$ & $-0.139$\\
就业率(\%) & $-0.492(-0.68)$ & $0.864(0.73)$ & $0.046$ & $-0.091$ & $0.005$\\
政府雇佣(\%) & $0.863(0.95)$ & $0.480(0.32)$ & $0.153$ & $-0.030$ & $0.007$\\
家庭收入中位数 & $0.009(0.16)$ & $0.015(0.17)$ & $-0.125$ & $0.082$ & $0.989$\\
方差 & $107.02$ & $39.67$ & $8.37$ & $2.87$ & $0.15$\\
累积总方差百分比 & $67.7$ & $92.8$ & $98.1$ & $99.9$ & $1.00$\\
\bottomrule
\end{tabular}
\end{mvtable}

第一主成分解释了样本总方差的 $67.7\%$. 前两个主成分合起来占样本总方差的 $92.8\%$. 因此, 样本变差可以用前两个主成分来很好地解释. 数据从关于 5 个变量的 61 个观测压缩到关于两个主成分的 61 个观测是合理的.

给定上述的主成分系数, 第一主成分本质上表示政府雇佣百分比对于就业总百分比的加权差. 第二主成分则表示两个变量的加权和. \hfill$\blacksquare$

正如我们在讨论总体主成分时说过的那样, 为了阐明主成分, 成分系数 $\hat{e}_{ik}$ 和相关系数 $r_{\hat{y}_1,x_k}$ 都应进行考查. 相关系数考虑到了原变量方差的不同, 但仅仅量测单个 $X$ 的重要性却未顾及组成该成分的其他 $X$. 然而, 从例 8.3 注意到, 表中的相关系数的确证实了由成分系数所提供的阐述.

\mvsub{主成分个数}

始终有个问题: 应保留多少个成分? 对此没有一个确切的回答. 应考虑的事情包括所说明的样本总方差的量、特征值(样本成分方差)的相对大小以及诸成分的现实题材的阐述. 此外, 像我们稍后讨论的那样, 与一个接近于零的特征值有关的成分, 会显得不重要, 却可能表示数据中有一个意外的线性依赖关系.

一种能帮助我们确定主成分合适个数的有用的视觉工具, 是所谓的崖底碎石图.\footnote{{\footnotesize 译者注: 此图得名于 $(i,\lambda_i)$ 连线的图形, 对应大特征值的点较高, 然后陡然下落成“悬崖”状, 对应后几个小特征值的点则似崖下一堆碎石.}} 将特征
